\section{Stopped interfaces and comparison of experiments}\label{sec:composition}
The lower and upper sides must refer to the same task under any experimental reduction. A complete-law defect is stronger than a predictive discrepancy and is added after root metric errors have been combined.

\subsection{Bounded stopped prediction}
For a visible stopping time $\tau\le mK_0$, define the stopped sigma field $\mathcal H_\tau$, a declared bounded task $Y^\tau$, and
\[
 P_\tau=\E[Y^\tau\mid\mathcal H_\tau],\qquad
 B_\tau=\E\norm{Y^\tau-P_\tau}^2.
\]
At $\tau=0$ these mean the actual preparation and its legal terminal task. Assume $P_\tau=f_\tau(S_\tau)$, with a common Lipschitz constant, and that every stopped exploration is the prefix of a specified legal continuation to $mK_0$ for which the block certificate holds. This continuation is used only for estimates; post-stop calls and information are not used by the actual predictor.

\begin{proposition}[A selection-sensitive stopped bound]\label{prop:stopped}
Under the upper hypotheses of Theorem~\ref{thm:transfer}, put $\rho=r_M+u$ and $\sigma=\lceil\tau/m\rceil$. There are explicit finite constants $c,C$, depending only on the block certificate, such that with
$V_0=d(S_0,\widehat S_0)^2+c\rho^2w(\widehat S_0)^2$,
\begin{equation}\label{eq:stoppedbound}
 \E d(S_\tau,\widehat S_\tau)^2
 \le C\{\E V_0+\rho^2\E\sigma\}.
\end{equation}
Consequently $e_0=O(\rho)$ and $\E\sigma\le L_0$ give
$R\le B_\tau+(C'\rho\sqrt{1+L_0}+v)^2$.
If the terminal decoder has a prediction label of size $M$ and a retained stopping clock of size $Q$, its checkpoint identity uses $q_{MQ}(\law(P_\tau))$ when this is the whole accessible tuple. A free declared interface instead requires conditioning on that interface.
\end{proposition}
\begin{proof}
Here are constants and the optional-sum argument, to distinguish the claim from substituting a random time into a deterministic estimate. Put
\[
 \bar\kappa=(1+\kappa)/2,\quad C_Y=\bar\kappa/(\bar\kappa-\kappa),
 \quad\gamma=\max\{\bar\kappa,(1+a)/2\},
\]
\[
 c=\max\{1,C_Y\Lambda_m^2A/(\gamma-a)\},\qquad
 c_0=C_Y\Lambda_m^2(B+1)+cb.
\]
The conditional shadow bound, $r_M\sqrt{x}+u\le\rho\sqrt{x+1}$ and
$(\sqrt\kappa d+z)^2\le\bar\kappa d^2+C_Yz^2$ give
\[
 \E[V_{k+1}\mid\mathcal H_{km}]\le\gamma V_k+c_0\rho^2,
 \quad V_k=d(S_{km},\widehat S_{km})^2+c\rho^2w(\widehat S_{km})^2.
\]
Multiply the rearranged inequality by the predictable indicator $\one_{\{k<\sigma\}}$ and sum the finitely many terms. This gives
$(1-\gamma)\E\sum_{k<\sigma}V_k\le\E V_0+c_0\rho^2\E\sigma$.
For a raw stopping time, the within-block telescoping bound and the inequality that a maximum of nonnegative terms is at most their sum give
\[
 \E[\max_{1\le j\le m}d(S_{km+j},\widehat S_{km+j})^2\mid\mathcal H_{km}]
 \le C_1V_k+C_2\rho^2,
\]
where $C_1=2m(L^{2m}+\Lambda_m^2A/c)$ and $C_2=2m\Lambda_m^2(B+1)$. Sum only over entered blocks; their indicator is known at each block start. At $\tau=0$ use $V_0$. The resulting bound is
\[
 (1+C_1/(1-\gamma))\E V_0+
 (C_1c_0/(1-\gamma)+C_2)\rho^2\E\sigma,
\]
which proves \eqref{eq:stoppedbound}. Stopped conditional projection proves the risk statement. The information-cut identity follows by counting every available terminal state before applying Lemma~\ref{lem:cut}; it is not an unqualified $M$-center formula when the clock carries information.
\end{proof}

\subsection{Measurable report coupling}
For probability kernels $P_z,Q_z$ on a standard Borel report space, a jointly measurable maximal coupling exists. One construction uses nested finite partitions generated by a countable separating algebra. Relative to $R_z=(P_z+Q_z)/2$, their atomwise density ratios are bounded measurable functions. Martingale convergence gives a jointly measurable clipped limsup version $p=dP_z/dR_z$ for each $z$; put $q=2-p$. The common kernel $\mu_z=\min(p,q)R_z$ has mass $c_z$. The coupling
\[
 \operatorname{diag}_*\mu_z+
 \frac{(P_z-\mu_z)\otimes(Q_z-\mu_z)}{1-c_z}\quad(c_z<1)
\]
and the diagonal law for $c_z=1$ are measurable kernels. The residual measures are mutually singular, so their product has zero diagonal mass. The disagreement probability is $1-c_z=\TV(P_z,Q_z)$, with $\TV=\sup_A|P(A)-Q(A)|$.

Suppose report kernels satisfy $\TV(J(\cdot\mid x,a),J(\cdot\mid z,a))\le L_Jd(x,z)$, and task kernels have constant $L_Y$. A preparation mismatch is either absent or explicitly at most $\varepsilon_0$. Successive maximal couplings, driven until first disagreement by the same declared controller and stopping rule, give complete stopped report-and-task defect at most
\begin{equation}\label{eq:reporttv}
 \min\{1,\varepsilon_0+L_J\E\sum_{t=1}^\tau d(S_{t-1},\widehat S_{t-1})
                         +L_Y\E d(S_\tau,\widehat S_\tau)\}.
\end{equation}
Indeed before disagreement both transcripts, actions and stopping decisions coincide. Continue a physical-law shadow after disagreement and sum first-disagreement probabilities, dropping the agreement indicator. This preserves the unconditional shadow estimates. A stopped sum is not replaced by $\E\tau$ times an unconditional deterministic-time bound. A finite complete-law bound is not infinite-path or large-deviation equivalence.

\subsection{Resource-certified morphisms}
For clarity we state all components of the comparison used here. A source experiment $\mathfrak E$ implements a target $\mathfrak F$ on a common parameter set by parameter-independent transducer initialization, state-update, source-action and report-transformation kernels. These kernels induce a strategy lift from the declared target class to legal source strategies. Each virtual completion time $\sigma_i$ is a source stopping time bounded by the given physical budget; for every permitted target stop $\tau$, $\sigma_\tau$ is a legal source stop and no later source information enters its decision. Legal actions and the available interface agree after each transformed prefix. Virtual costs, preparation and failures are included.

There is one Borel common scored-outcome map and one decision space, with loss in $[0,L_*]$. Uniformly over the common parameters, declared target policies and permitted bounded stops, the joint law of virtual history, permitted public policy coins and scored outcome differs from its target law by at most $\varepsilon$ in total variation. A transducer relying on a real hidden history or unspecified phase is not a finite implementation. Let its state count be $R$, target prediction count $M$, controller count $C_{\rm ctr}$ and phase/clock count $Q$. Workspace, read-only program, precision, raw calls and physical time are separate coordinates. The companion~\cite{companion} gives the fully typed standard Borel transcript construction; these requirements are precisely the ones needed in the following implication.

\begin{theorem}[Common-risk composition and reverse transfer]\label{thm:morphism}
Such a certificate implements a target risk $r_\lambda$ by a serial source tuple of at most $MR C_{\rm ctr}Q$ states and risk at most $r_\lambda+L_*\varepsilon$. Defects of two composable certificates add, capped at one, provided the lifted policies, stopping rules and common-task maps lie in the next certified class. Serial state overheads multiply; raw costs sum over executed subroutines, and nested per-call upper bounds may multiply.

If every source $K$-state procedure can conversely be implemented in the target with at most $\Psi(K)$ states and complete defect $\varepsilon_-$, then for the same task and risk criterion
\begin{equation}\label{eq:morphismlower}
 R_E(K)\ge R_F(\Psi(K))-L_*\varepsilon_-.
\end{equation}
In particular every acquired continuation lower of Lemma~\ref{lem:cut} transfers with that inflated budget and the target baseline left unchanged. A block upper transfers as
\begin{equation}\label{eq:morphismupper}
 R_E(MRC_{\rm ctr}Q)
 \le B_T+(L_fE_{M,u}+v)^2+L_*\varepsilon.
\end{equation}
The same assertions hold for minimax risks when both certificates are uniform over the same parameter set.
\end{theorem}
\begin{proof}
Run the physical subroutine and retain the tuple of predictor, simulator, controller and clock states. The action and stopping conditions make the procedure executable; its cardinality is the displayed product. Applying the common decision kernel preserves total-variation closeness. The bounded-loss inequality for our convention gives a difference at most $L_*\varepsilon$.

For composition insert the intermediate stopped joint law and use the triangle inequality and contraction of total variation under the downstream parameter-independent transformation. Class closure is needed here for adaptive lifted policies. Retaining both serial tuples proves the storage upper. All actual calls, waiting times and failures are summed; shared program or workspace can sometimes be reused, so additive serial descriptions are only upper accounts.

For the reverse statement apply the upper implication to every source procedure: $R_F(\Psi(K))\le r_E+L_*\varepsilon_-$. Infimize over source procedures. Uniformity permits the same argument after taking the parameter supremum. No separately recomputed source and target Bayes baselines are subtracted. This proves both \eqref{eq:morphismlower} and \eqref{eq:morphismupper} in their stated directions.
\end{proof}

The exact-bit and erasure pair in Theorem~\ref{thm:joint} verifies this definition with a prescribed completion clock, parameter-independent bit transformation and a common terminal audit; it is not just a marginal report comparison. The Gaussian report comparison verifies \eqref{eq:reporttv} for the same declared controller. These are distinct uses: the latter does not manufacture a reverse certificate for all Gaussian source algorithms. Likewise a product $RC_{\rm ctr}Q$ is not necessary merely because one implementation stores it. Only an executable challenge separating those registers, or an appropriate reverse theorem, provides that lower bound.
